\section{Two-sided physical ends and oscillation of the original source}
\label{sec:v90-two-sided}
Fix an exact event, and use its Borel maximal-segment coordinates
$(\sigma,u)$, with $a_\sigma<u<b_\sigma$. The measure $d\sigma$
includes the disjoint angular charts and their incident components.
Write $A_\sigma(u)$ for the physical density, $z=|g|$, and
$\ell_\sigma=b_\sigma-a_\sigma$. In this section
\[
 \gamma=\frac1{32},\qquad \eta=\frac14,\qquad
 \kappa=\frac{1-\eta}{2}=\frac38,\qquad v=v^{1/4}.
\]
Every estimate is on the original collision source. In particular a
backward end is an end of the roof flow in the same exact event, not
a reversal or relabeling of the billiard record.

\begin{lemma}[Both physical end layers]\label{lem:v90-ends}
Let
\[
 e_s^\pm(t)=\int\one_{\{\tau^\pm\le s\}}\,d\mu_t.
\]
Backward critical termination is included in $e_s^-$. There are
$s_0,C>0$ such that
\begin{equation}\label{eq:v90-end-cap}
 \bar\nu\{\tau^\pm\le s\}\le C\sqrt s\qquad(0<s<s_0).
\end{equation}
For every fixed auxiliary band $B_*$, after increasing its constants
and choosing one $0<\theta_{B_*}<1$, both signs satisfy
\begin{align}
 m^2\int_J e_s^\pm(t)\,dt
 \le{}& Cs^\gamma(|J|+B_*^{-1})\notag\\
 &+C_{B_*}(1+B_*|J|)m^3
                        \min\{\theta_{B_*}^{m},\sqrt s\}.
                                             \label{eq:v90-end-finite}
\end{align}
The constants are uniform in all finite $m\ge2$, radii and exact labels.
\end{lemma}
\begin{proof}
The forward inclusion is \eqref{eq:v83-layer-input}. For a backward
segment of length at most $s$, apply the same contact-motion estimate
to the segment with its earlier endpoint as the starting point. Its
endpoint-curve length is at most $C\sqrt s$, and each contact and
flight margin changes by at most its prescribed depth weight times
that length. If the limiting backward end is regular and noncritical,
strictly positive incidences, clearances and section-edge distances
would continue both the word and the field, contradicting maximality.
Thus the same three physical layers contain this case. If the limit
is critical, radial growth gives instead
\[
                     z(x)^2\le2\Lambda_\circ s.
\]
These alternatives also cover simultaneous failures. No chart seam
is used as an end.

The physical estimates used in the complete residence cap apply to
each of the three layers: incidence has mass at most $w^2$, a section
strip has mass at most $Cw$, and the image-collision clearance envelope
has mass at most $Cw$. At widths $Cq_0^d\sqrt s$ their complete depth
sum is $O(\sqrt s)$. The additional critical set lies in
$\{|p_0|\le C\sqrt s\}$, of invariant mass $O(\sqrt s)$. Dropping
the section tests on these positive upper comparisons costs exactly
$c_*^{-1}$. This proves \eqref{eq:v90-end-cap} for the original
source, without changing its labels.

For the localized estimate, the incidence and clearance layers are
contained in the original bad source at an auxiliary parameter
$C\sqrt s$. The section layers use \eqref{eq:v83-strip-input},
summed at their depth widths. This is the forward finite-count exit
proof and applies to the same layers on the backward side. For the
critical alternative use the retained endpoint-only bound
\[
 m^2\bar\nu\{F\in J,z<\zeta\}
 \le C\zeta^{129/128}(|J|+B_*^{-1})
       +C_{B_*}(1+B_*|J|)m^2\zeta\rho_{B_*}^{m/2}.
\]
At $\zeta=C\sqrt s$ its leading term is smaller than $Cs^\gamma$.
This proves the uncapped version of \eqref{eq:v90-end-finite}.
Apply $x\le P+D$, $x\le V\Rightarrow x\le P+\min(D,V)$ to
that positive mass and \eqref{eq:v90-end-cap}. Since $m^2\le m^3$,
adjusting constants gives the displayed common cap. It is a cap on
the positive estimate, not a claim that an isolated signed spectral
remainder has the physical bound.
\end{proof}

For use below the complete sources $p$ and $v$ obey the retained
finite-count estimate
\begin{equation}\label{eq:v90-volume-finite}
 m^2\int_J(p+v)\le U_{m,B_*}(|J|),\quad
 U_{m,B_*}(L)=C(L+B_*^{-1})
                   +C_{B_*}(1+B_*L)m^3\theta_{B_*}^{m}.
\end{equation}
Its constants may be enlarged from those in
\eqref{eq:v90-end-finite}. The stated fractional source, not an
untruncated first inverse gradient, is used here.

\begin{lemma}[A fixed-band envelope at every microscopic scale]
\label{lem:v90-remainder-envelope}
For $0<\theta<1$ and $0<s<1$,
\begin{equation}\label{eq:v90-cap-log}
 \sup_{m\ge2}m^3\min\{\theta^m,\sqrt s\}
          \le C_\theta\sqrt s(1+|\log s|)^3.
\end{equation}
Consequently, at every one fixed admissible band and fixed $L$,
\begin{equation}\label{eq:v90-uniform-end}
 \sup_{m,R,n,k}\sup_{|J|\le L}m^2\int_J e_s^\pm
                \le K_{B_*,L}s^\gamma,\qquad
 \sup_m U_{m,B_*}(L)<\infty.
\end{equation}
These suprema precede $s\downarrow0$.
\end{lemma}
\begin{proof}
Put $a=-\log\theta$ and $M_s=\max\{2,a^{-1}\log(s^{-1/2})\}$.
For $m\le M_s$, use $m^3\sqrt s$. For $m>M_s$, write
$m=M_s+y$ and use
$(M_s+y)^3e^{-ay}\le C_a(1+M_s)^3$; when $M_s=2$ the
same estimate follows after increasing $C_a$. This proves
\eqref{eq:v90-cap-log}. Since $\gamma<1/2$,
$s^{1/2-\gamma}(1+|\log s|)^3$ is bounded on $(0,1)$.
Substitution in \eqref{eq:v90-end-finite} proves the first assertion
of \eqref{eq:v90-uniform-end}. Exponential decay proves the second.
No Fourier band is allowed to grow with the collision count.
\end{proof}

Extend
\[
 f_\sigma^\varepsilon(u)=A_\sigma(u)W^\varepsilon(\Psi_\sigma(u))
\]
by zero outside $(a_\sigma,b_\sigma)$. Define the lifted oscillation
\begin{equation}\label{eq:v90-lifted-oscillation}
 \Xi_h^\varepsilon(J)=\int_J\int
    |f_\sigma^\varepsilon(t+h)-f_\sigma^\varepsilon(t)|\,d\sigma\,dt.
\end{equation}
This keeps the absolute value before the complete physical sum.
It therefore dominates $\int_J|b^\varepsilon(t+h)-b^\varepsilon(t)|dt$.

\begin{theorem}[All-scale oscillation of the original weighted source]
\label{thm:v90-whole-oscillation}
For each fixed $L$ and one fixed admissible band there is $K_{B_*,L}$
such that, for every $m\ge2$, every exact event, every $0<h<h_0$ and
$0<\varepsilon\le1$,
\begin{equation}\label{eq:v90-whole-oscillation}
 m^2\Xi_h^\varepsilon(J)
       \le K_{B_*,L}\bigl(h^{1/32}+\varepsilon^{-1}\sqrt h\bigr)
                         \qquad(|J|\le L).
\end{equation}
The estimate includes every true forward end, backward end and
critical entrance. It is not restricted to common survivors.
For $p$, the unweighted version has no guard-derivative term.
\end{theorem}
\begin{proof}
The all-contact first derivatives and the product rule for factors in
$[0,1]$ give
\[
 |D_qW^\varepsilon|\le
 C\varepsilon^{-1}\sum_{d\ge0}q_0^{3d/4}
 \le C_\Theta\varepsilon^{-1}.
\]
This uses the original smooth guard and requires no monotonicity of
$\Theta$. Since $|X_q|\le C/z$, its derivative on each regular
segment satisfies $|(W^\varepsilon)'|\le C_\Theta\varepsilon^{-1}/z$
almost everywhere. Together with the full angular-density derivative,
\begin{equation}\label{eq:v90-weighted-derivative}
 |(f_\sigma^\varepsilon)'|
       \le C A_\sigma I_m/z+C_\Theta\varepsilon^{-1}A_\sigma/z.
\end{equation}
Local absolute continuity suffices; all subsequent integrals are
first taken on compact subsegments.

For a pair $t,t+h$ on the same segment, integrate
\eqref{eq:v90-weighted-derivative}. In the subsequent integral over
$t\in J$, an intermediate roof $u$ is counted for a set of starting
roofs of length at most $\min\{h,u-a_\sigma\}$. Radial growth from
the lower endpoint gives $z(u)^2\ge2\lambda_\circ(u-a_\sigma)$.
Thus
\begin{align*}
 \frac{\min(h,u-a_\sigma)}{z(u)}
     &\le C h^{5/8}z(u)^{-1/4},\\
 \frac{\min(h,u-a_\sigma)}{z(u)}&\le C\sqrt h.
\end{align*}
Use the first inequality for optical variation and the second for
the guard. By the Borel bijection their complete sums are bounded by
\[
 C h^{5/8}\int_{J+(0,h)}v(u)\,du
        +C_\Theta\varepsilon^{-1}\sqrt h\int_{J+(0,h)}p(u)\,du.
\]
There is no physical-word multiplicity. Tonelli and positive
exhaustion justify the sum even when the unweighted derivative
$A_\sigma I_m/z$ is not globally integrable.

If only $t$ is on a segment, its contribution is at most $e_h^+(t)$.
If only $t+h$ is on it, its contribution is at most $e_h^-(t+h)$.
These are the actual original end layers. Therefore the preceding
drift bound plus $\int_Je_h^++\int_{J+h}e_h^-$ bounds
$\Xi_h^\varepsilon(J)$. Apply \eqref{eq:v90-volume-finite} and
\eqref{eq:v90-uniform-end}. Since $h^{5/8}\le h^{1/32}$ for
$h<1$, the result follows. Negative translations follow by translating
$J$ and interchanging its two arguments. The common segment kernels
fix the same representatives throughout; this is an integral identity
for all shifts, not a supremum over exceptional fiber versions.
\end{proof}
